\chapter{Esercizi: Realizzazione di classi in Java}
\label{chap:esercizi-realizzazione-classi-java}

\begin{exercise}
\label{ex:fi-esercizi-java-classi-modello}
Che cosa significa dire che una classe è un ``modello'' per creare oggetti?
\end{exercise}

Una classe definisce la struttura e il comportamento condivisi dalle
sue istanze, specificando campi, costruttori e metodi. Ogni oggetto
creato secondo tale modello può possedere un proprio stato
(\cref{sec:fi-java-classi-stato}).

\begin{exercise}
\label{ex:fi-esercizi-java-classi-classe-oggetto}
Qual è la differenza tra una classe e un oggetto?
\end{exercise}

La classe è la definizione di un tipo, mentre un oggetto è una sua
istanza concreta, creata durante l'esecuzione. Per esempio,
\texttt{BankAccount} è la classe, mentre
\texttt{new BankAccount()} crea un oggetto di quel tipo
(\cref{sec:fi-java-classi-stato,sec:fi-java-classi-costruttori}).

\begin{exercise}
\label{ex:fi-esercizi-java-classi-collaudo-scopo}
A che cosa serve un programma di collaudo per una classe Java?
\end{exercise}

Serve a verificarne il comportamento mediante la costruzione di
oggetti e l'invocazione dei loro metodi, confrontando i risultati
ottenuti con quelli attesi. Può individuare errori, ma il superamento
di alcuni test non dimostra la correttezza generale
(\cref{sec:fi-java-collaudo,sec:fi-java-classi-collaudo}).

\begin{exercise}
\label{ex:fi-esercizi-java-classi-collaudo-passi}
Quali sono i passaggi principali per costruire un programma di collaudo?
\end{exercise}

Si definisce una classe di collaudo con un metodo \texttt{main},
si creano istanze della classe da verificare, se ne invocano i
metodi con ingressi appropriati e si confrontano gli stati o i
valori ottenuti con i risultati previsti, includendo i casi limite
(\cref{sec:fi-java-classi-collaudo}).

\begin{exercise}
\label{ex:fi-esercizi-java-classi-black-box}
Che cosa si intende per approccio \emph{black-box} nella programmazione
orientata agli oggetti?
\end{exercise}

Consiste nell'utilizzare un oggetto attraverso la sua interfaccia
pubblica e il comportamento dichiarato, senza dover conoscere i
dettagli della sua implementazione interna
(\cref{sec:fi-java-classi-black-box}).

\begin{exercise}
\label{ex:fi-esercizi-java-classi-interfaccia}
Che cosa si intende per interfaccia pubblica di una classe?
\end{exercise}

È l'insieme dei membri accessibili ai suoi utilizzatori, con le
modalità e le condizioni del loro impiego. Comprende, in particolare,
i costruttori e i metodi pubblici, i relativi parametri, i risultati
e il comportamento previsto (\cref{def:fi-java-classi-interfaccia}).

\begin{exercise}
\label{ex:fi-esercizi-java-classi-intestazione-metodo}
Da quali parti è composta l'intestazione di un metodo Java?
\end{exercise}

L'intestazione comprende gli eventuali modificatori, il tipo
restituito (\texttt{void} se non restituisce un valore), il nome
del metodo e l'elenco dei parametri, ciascuno dichiarato con tipo
e nome. Il corpo del metodo, distinto dall'intestazione, è racchiuso
fra parentesi graffe (\cref{sec:fi-java-classi-firma}).

\begin{exercise}
\label{ex:fi-esercizi-java-classi-accesso-modifica}
Che differenza c'è tra un metodo di accesso e un metodo modificatore?
\end{exercise}

Un metodo di accesso osserva lo stato di un oggetto senza modificarlo;
un metodo modificatore può invece alterarlo
(\cref{def:fi-java-classi-accessori}).

\begin{exercise}
\label{ex:fi-esercizi-java-classi-get-balance}
Perché \texttt{getBalance()} è un metodo di accesso?
\end{exercise}

Restituisce il saldo memorizzato nell'oggetto
\texttt{BankAccount} senza modificarne il campo
\texttt{balance} (\cref{def:fi-java-classi-accessori,sec:fi-java-classi-incapsulamento}).

\begin{exercise}
\label{ex:fi-esercizi-java-classi-deposit-withdraw}
Perché \texttt{deposit()} e \texttt{withdraw()} sono metodi modificatori?
\end{exercise}

Entrambi cambiano il campo \texttt{balance}: il primo aumenta
il saldo dell'importo versato e il secondo lo diminuisce
dell'importo prelevato
(\cref{def:fi-java-classi-accessori,sec:fi-java-classi-incapsulamento}).

\begin{exercise}
\label{ex:fi-esercizi-java-classi-costruttore}
Che cos'è un costruttore in Java e qual è il suo scopo principale?
\end{exercise}

È una dichiarazione speciale eseguita durante la creazione di
un oggetto per inizializzarne lo stato. Il costruttore viene
richiamato durante la valutazione di un'espressione con
\texttt{new} (\cref{def:fi-java-classi-costruttore}).

\begin{exercise}
\label{ex:fi-esercizi-java-classi-costruttore-differenze}
Quali due caratteristiche distinguono un costruttore da un normale metodo?
\end{exercise}

Il costruttore ha necessariamente lo stesso nome della classe e
non dichiara un tipo di ritorno, neppure \texttt{void}
(\cref{def:fi-java-classi-costruttore}).

\begin{exercise}
\label{ex:fi-esercizi-java-classi-costruttore-void}
Perché il costruttore \texttt{public BankAccount()} non deve dichiarare
\texttt{void}?
\end{exercise}

Un costruttore non ha un tipo restituito: scrivendo
\texttt{public void BankAccount()} si definirebbe invece
un metodo ordinario chiamato \texttt{BankAccount}, non un
costruttore (\cref{sec:fi-java-classi-costruttori}).

\begin{exercise}
\label{ex:fi-esercizi-java-classi-piu-costruttori}
Che cosa significa avere più costruttori nella stessa classe?
\end{exercise}

Significa offrire diverse modalità di inizializzazione.
Il costruttore senza parametri può creare un conto con saldo
nullo; quello che riceve un parametro \texttt{double}
può utilizzare il saldo iniziale indicato dal chiamante
(\cref{sec:fi-java-classi-costruttori}).

\begin{exercise}
\label{ex:fi-esercizi-java-classi-overloading}
Che cos'è l'\emph{overloading} e come fa il compilatore a scegliere
quale costruttore invocare?
\end{exercise}

L'\emph{overloading} consente di dichiarare più costruttori o
metodi con lo stesso nome ma parametri differenti per numero,
tipo od ordine. Il compilatore seleziona la dichiarazione
applicabile agli argomenti effettivi secondo le regole delle
conversioni e della specificità; un'invocazione non applicabile
o ambigua produce un errore di compilazione
(\cref{def:fi-java-classi-overloading}).

\begin{exercise}
\label{ex:fi-esercizi-java-classi-private-balance}
Perché la variabile di istanza \texttt{balance} viene dichiarata
\texttt{private} nella classe \texttt{BankAccount}?
\end{exercise}

Per impedire modifiche dirette dall'esterno e consentire l'accesso
al saldo attraverso i metodi pubblici previsti. Questa separazione
realizza l'incapsulamento e permette di controllare la validità
dello stato, purché i costruttori e i metodi modificatori effettuino
le verifiche necessarie
(\cref{sec:fi-java-classi-incapsulamento,warn:fi-java-classi-bank-limiti,sec:fi-java-cicli-shadowing}).
